%% GENERATED by tools/gen_latex.py -- edit content/ch05.py instead.
%% Build:  latexmk -pdf ch05.tex     (or: pdflatex ch05.tex, twice)
\documentclass[aspectratio=169,11pt,t]{beamer}
\usepackage{itbpro}

\coursetitle{Dasar-Dasar Machine Learning}
\coursesubtitle{Ketegangan antara optimalisasi dan generalisasi -- mengapa ia tak terhindarkan, bagaimana mengukurnya, dan apa saja yang benar-benar menolong.}
\coursekicker{Chapter 5}
\coursesource{Chollet \& Watson, Deep Learning with Python 3e -- bab 5 (hlm. 136-170)}
\courseduration{3 jam (2 sesi)}
\coursepresenter{Prof. Bambang Riyanto Trilaksono}{Pengajar Utama}
\coursebrand{AI for Professional \textperiodcentered\ ITB \textperiodcentered\ Ch 5}

\title{Dasar-Dasar Machine Learning}
\author{Prof. Bambang Riyanto Trilaksono}
\date{}

\begin{document}
\covertitle

\framekicker{Learning outcomes}
\begin{frame}[shrink=25]{Tujuan Sesi}
\leadpar{Setelah sesi ini, peserta dapat:}
\begin{isteps}
\item Merumuskan \textbf{ketegangan optimalisasi vs generalisasi}, dan menyebut tiga penyebab overfitting: data berderau, fitur ambigu, dan fitur langka.
\item Menjelaskan \textbf{hipotesis manifold} dan mengapa generalisasi deep learning adalah \textit{interpolasi}, bukan penalaran.
\item Memilih protokol evaluasi yang tepat -- \textbf{hold-out, K-lipat, K-lipat berulang} -- dan menghindari tiga jebakannya.
\item Menetapkan \textbf{tolok banding akal sehat} sebelum melatih apa pun.
\item Mendiagnosis tiga kegagalan pelatihan: \textbf{tidak mulai, tidak menggeneralisasi, tidak bisa overfit} -- dan tahu obat masing-masing.
\item Menerapkan \textbf{kurasi data, rekayasa fitur, early stopping, pengurangan kapasitas, regularisasi bobot, dan dropout}, serta tahu kapan masing-masing yang tepat.
\end{isteps}

\vskip 6pt
\vskip 4pt
\reslink{Buku}{Bab 5 --- teks penuh}{https://deeplearningwithpython.io/chapters/chapter05\_fundamentals-of-ml}
\reslink{Notebook}{01\_korelasi\_semu\_dan\_derau.ipynb}{../../notebooks/ch05/01\_korelasi\_semu\_dan\_derau.ipynb}
\reslink{Notebook}{02\_label\_diacak\_manifold.ipynb}{../../notebooks/ch05/02\_label\_diacak\_manifold.ipynb}
\reslink{Notebook}{03\_protokol\_evaluasi.ipynb}{../../notebooks/ch05/03\_protokol\_evaluasi.ipynb}
\reslink{Notebook}{04\_kapasitas\_model.ipynb}{../../notebooks/ch05/04\_kapasitas\_model.ipynb}
\reslink{Notebook}{05\_regularisasi\_l2\_dropout.ipynb}{../../notebooks/ch05/05\_regularisasi\_l2\_dropout.ipynb}
\reslink{Repo}{Notebook resmi penulis}{https://github.com/fchollet/deep-learning-with-python-notebooks/blob/master/chapter05\_fundamentals-of-ml.ipynb}
\vskip 2pt
\end{frame}

\framekicker{Peta bab}
\begin{frame}[shrink=25]{Dari 'modelnya jalan' ke 'modelnya boleh dipakai'}
\leadpar{Bab 4 memperkenalkan overfitting sebagai kejadian. Bab 5 menjadikannya \hilite{kerangka berpikir}, dan hampir semua praktik baik di sisa buku ini adalah cara menangani satu ketegangan yang sama.}
\begin{minipage}[t]{0.241\textwidth}
\begin{icardaccent}
\cardh{5.1 $\cdot$ Generalisasi}
Apa yang menyebabkan overfitting, dan mengapa deep learning bisa menggeneralisasi sama sekali.
\end{icardaccent}
\end{minipage}%
\hfill
\begin{minipage}[t]{0.241\textwidth}
\begin{icardaccent}
\cardh{5.2 $\cdot$ Menilai model}
Tiga protokol, tolok banding akal sehat, dan tiga jebakan yang membatalkan penilaian.
\end{icardaccent}
\end{minipage}%
\hfill
\begin{minipage}[t]{0.241\textwidth}
\begin{icardaccent}
\cardh{5.3 $\cdot$ Memperbaiki fit}
Tiga kegagalan pelatihan dan obatnya. Sasarannya: \textbf{bisa} overfit dulu.
\end{icardaccent}
\end{minipage}%
\hfill
\begin{minipage}[t]{0.241\textwidth}
\begin{icardaccent}
\cardh{5.4 $\cdot$ Memperbaiki generalisasi}
Kurasi data, rekayasa fitur, early stopping, dan tiga cara regularisasi.
\end{icardaccent}
\end{minipage}%
\begin{iquote}
\quotetext{Persoalan pokok dalam machine learning adalah ketegangan antara \textbf{optimalisasi} dan \textbf{generalisasi}. Tujuannya jelas generalisasi yang baik -- tetapi Anda tidak mengendalikan generalisasi; yang bisa Anda lakukan hanya memasangkan model ke data latihnya.}\quotecite{Chollet \& Watson, bab 5.1}
\end{iquote}
\end{frame}
\note{Kalimat kuncinya: 'Anda tidak mengendalikan generalisasi.' Semua teknik di bab ini adalah cara tidak langsung --- tak satu pun mengatur generalisasi secara langsung.}

\coursesection{01}{Generalisasi: tujuan machine learning}{Mengapa overfitting terjadi di setiap persoalan, tanpa kecuali.}

\framekicker{Bagian 5.1.1}
\begin{frame}[shrink=25]{Kurva kanonis -- dan ia berlaku universal}
\begin{center}
\adjustbox{max width=\linewidth,max totalheight=0.52\textheight}{%

\begin{tikzpicture}[font=\sffamily\tiny]
  \fill[itbbluelt!8]  (0,0) rectangle (2.6,3.2);
  \fill[limebr!10]    (2.6,0) rectangle (4.0,3.2);
  \fill[rosebr!9]     (4.0,0) rectangle (9.6,3.2);
  \draw[rule, line width=0.8pt] (0,0) -- (9.6,0);
  \draw[rule, line width=0.8pt] (0,0) -- (0,3.2);
  \node[text=ink3, anchor=north] at (4.8,-0.15) {waktu pelatihan};
  \node[text=ink3, rotate=90, anchor=south] at (-0.35,1.6) {rugi};
  \node[text=itbblue] at (1.3,3.0) {underfitting};
  \node[text=lime]    at (3.3,3.0) {robust fit};
  \node[text=rose]    at (6.6,3.0) {overfitting};
  \draw[signal, line width=1.2pt]
    (0,2.75) .. controls (1.6,1.75) and (3.0,1.1) .. (5.0,0.72)
             .. controls (7.0,0.42) and (8.4,0.26) .. (9.6,0.2);
  \draw[rose, line width=1.2pt] (0,2.6) .. controls (1.4,1.7) and (2.4,1.28) .. (3.3,1.18);
  \draw[rose, line width=1.2pt] (3.3,1.18) .. controls (5.2,1.3) and (7.6,2.1) .. (9.6,3.05);
  \draw[amberbr, line width=0.9pt, dashed] (3.3,0) -- (3.3,3.2);
  \fill[amberbr] (3.3,1.18) circle (2.2pt);
  \node[text=amber, anchor=west] at (0,-0.6)
    {Gambar 5.1 --- pola ini muncul pada SETIAP jenis model dan SETIAP kumpulan data};
\end{tikzpicture}

}
\end{center}
\vskip 3pt{\color{ink3}\fontsize{7.4}{9.4}\selectfont Gambar 5.1 -- di awal, optimalisasi dan generalisasi berjalan searah. Setelah sekian iterasi, keduanya berpisah.\par}
\begin{ibullets}
\item \textbf{Underfit} -- masih ada kemajuan yang bisa diraih; model belum menangkap semua pola yang relevan.
\item \textbf{Robust fit} -- titik terbaik. Sempit, dan hanya terlihat lewat metrik validasi.
\item \textbf{Overfit} -- model mulai mempelajari pola yang \hilite{khas data latih} tetapi menyesatkan pada data baru.
\end{ibullets}
\end{frame}

\framekicker{Bagian 5.1.1}
\begin{frame}[shrink=25]{Tiga penyebab overfitting}
\begin{minipage}[t]{0.3253\textwidth}
\begin{icardwarn}
\cardh{1 $\cdot$ Data berderau}
Masukan tak sah (citra MNIST hitam seluruhnya), dan yang lebih buruk: masukan sah yang \textbf{salah label}. Model yang memaksa mencakup pencilan ini akan salah pada data mirip.
\end{icardwarn}
\end{minipage}%
\hfill
\begin{minipage}[t]{0.3253\textwidth}
\begin{icardwarn}
\cardh{2 $\cdot$ Fitur ambigu}
Data bersih pun bisa berderau bila persoalannya memang tak pasti. Pisang \textit{mentah} vs \textit{matang} tak punya batas objektif; tekanan udara yang sama kadang diikuti hujan, kadang tidak.
\end{icardwarn}
\end{minipage}%
\hfill
\begin{minipage}[t]{0.3253\textwidth}
\begin{icardbad}
\cardh{3 $\cdot$ Fitur langka \& korelasi semu}
Kata \textit{cherimoya} muncul sekali, kebetulan di ulasan negatif $\rightarrow$ model memberinya bobot besar. \textbf{Tidak perlu langka-langka amat}: kata yang muncul 100 kali dengan 54\% positif sudah cukup.
\end{icardbad}
\end{minipage}%
\begin{iband}[rose]
Selisih 54\% lawan 46\% itu bisa saja \hilite{kebetulan statistik murni}. Model tetap akan memakainya. Chollet menyebut ini \textbf{salah satu sumber overfitting yang paling umum}.
\end{iband}
\end{frame}
\note{Contoh setara di data operasional: kode cabang atau jam transaksi yang kebetulan berkorelasi dengan label di data historis.}

\framekicker{Bagian 5.1.1 $\cdot$ listing 5.1-5.3}
\begin{frame}[shrink=25]{Percobaan: tambahkan derau murni, akurasi turun}
\icode{python}{listing 5.1 --- dua kumpulan pembanding}{listings/ch05-01.py}
\begin{minipage}[t]{0.575\textwidth}
{\fontsize{9.4}{12.6}\selectfont\color{fg2}Kandungan informasinya \textbf{identik} pada kedua kumpulan. Manusia tidak akan terpengaruh sama sekali oleh penambahan ini.\par}\vskip 4pt
\begin{iband}[rose]
Tetapi akurasi validasi model yang dilatih dengan kanal derau berakhir \hilite{sekitar satu poin persen lebih rendah} -- murni akibat korelasi semu. Makin banyak kanal derau, makin jauh merosotnya.
\end{iband}

\end{minipage}%
\hfill
\begin{minipage}[t]{0.395\textwidth}
\begin{minipage}[t]{1.0\textwidth}
\istat{$-$1 poin}{akurasi validasi, hanya karena derau ditempelkan}
\end{minipage}%

\end{minipage}%
{\fontsize{9.4}{12.6}\selectfont\color{fg2}Obatnya: \textbf{seleksi fitur}. Hitung skor kegunaan tiap fitur -- misalnya \textit{mutual information} antara fitur dan label -- lalu simpan yang di atas ambang. Membatasi IMDB ke 10.000 kata tersering di bab 4 adalah bentuk kasarnya.\par}\vskip 4pt
\end{frame}

\framekicker{Bagian 5.1.2 $\cdot$ listing 5.4}
\begin{frame}[shrink=25]{Model bisa memasangkan diri ke APA SAJA}
\icode{python}{listing 5.4 --- label diacak}{listings/ch05-02.py}
\ioutput{listings/ch05-03.txt}
\begin{iband}[amber]
Rugi latih tetap turun walau \textbf{tidak ada hubungan apa pun} antara masukan dan label. Model hanya \hilite{menghafal, seperti kamus Python}. Jadi: kemampuan memasangkan diri bukan bukti bahwa persoalannya bisa diselesaikan.
\end{iband}
{\fontsize{9.4}{12.6}\selectfont\color{fg2}Kalau begitu, mengapa deep learning menggeneralisasi sama sekali? Jawabannya, kata Chollet, \textbf{sedikit sekali berkaitan dengan modelnya}, dan banyak berkaitan dengan struktur informasi di dunia nyata.\par}\vskip 4pt
\end{frame}
\note{Ini slide diagnostik yang paling berguna: kalau rugi latih turun tapi validasi mandek di tolok banding, persoalannya bukan model --- datanya memang tidak memuat jawabannya.}

\framekicker{Bagian 5.1.2}
\begin{frame}[shrink=25]{Hipotesis manifold}
\begin{center}
\adjustbox{max width=\linewidth,max totalheight=0.52\textheight}{%

\begin{tikzpicture}[font=\sffamily\tiny]
  \node[draw=rule, fill=papertint, rounded corners=4pt, minimum width=4.6cm,
        minimum height=2.5cm] (p) at (0,0) {};
  \node[text=ink3, anchor=north west] at ($(p.north west)+(0.15,-0.1)$)
    {ruang induk (semua larik $28\times28$)};
  \draw[signal, line width=1.1pt] (-1.7,-0.75) .. controls (-0.8,0.5) and (0.8,0.5) .. (1.75,-0.55);
  \fill[signal] (-1.55,-0.5) circle (2.4pt);
  \fill[signal] (1.6,-0.42) circle (2.4pt);
  \fill[lime]   (0.0,0.22) circle (2.0pt);
  \draw[rose, line width=0.9pt, dashed] (-1.55,-0.5) -- (1.6,-0.42);
  \fill[rose]   (0.0,-0.46) circle (2.0pt);
  \node[text=lime, anchor=south] at (0,0.32) {angka sah};
  \node[text=rose, anchor=north] at (0,-0.56) {bukan angka};
  \node[text=signal, anchor=south west] at ($(p.south west)+(0.15,0.1)$) {manifold laten};

  \node[draw=rule, fill=papertint, rounded corners=4pt, minimum width=4.8cm,
        minimum height=2.5cm, anchor=west] (q) at (2.7,0) {};
  \node[anchor=north west, align=left, text=ink2] at ($(q.north west)+(0.18,-0.12)$)
    {\textcolor{signal}{\bfseries interpolasi pada manifold}\\
     titik antara tetap angka yang sah\\[3pt]
     \textcolor{rose}{\bfseries interpolasi linear}\\
     rerata piksel dua angka\\
     biasanya BUKAN angka};
  \node[text=amber, anchor=south west] at ($(q.south west)+(0.18,0.1)$)
    {Gambar 5.8 --- keduanya sama sekali berbeda};

  \node[text=ink3, anchor=north west, align=left] at (-2.4,-1.45)
    {$256^{784}$ larik mungkin --- jauh lebih banyak dari jumlah atom di alam semesta.\\
     Yang benar-benar berupa tulisan tangan hanya menempati sepetak kecil, dan petak itu bersinambung.};
\end{tikzpicture}

}
\end{center}
\vskip 3pt{\color{ink3}\fontsize{7.4}{9.4}\selectfont Gambar 5.7-5.8 -- angka tulisan tangan membentuk manifold di dalam ruang semua larik 28$\times$28 yang mungkin.\par}
\begin{iquote}
\quotetext{Hipotesis manifold menyatakan bahwa semua data alami terletak pada manifold berdimensi rendah di dalam ruang berdimensi tinggi tempat ia disandikan.}\quotecite{Chollet \& Watson, bab 5.1.2}
\end{iquote}
\begin{ibullets}
\item \textbf{Manifold} = subruang berdimensi lebih rendah yang secara setempat menyerupai ruang linear. Kurva mulus di bidang adalah manifold 1D dalam ruang 2D.
\item Ia \textbf{bersinambung}: ubah satu sampel sedikit, ia tetap angka yang sama. Dua angka mana pun dihubungkan oleh \hilite{lintasan mulus}.
\item Berlaku untuk MNIST, wajah manusia, bentuk pohon, suara manusia, bahkan bahasa alami.
\end{ibullets}
\end{frame}

\framekicker{Bagian 5.1.2}
\begin{frame}[shrink=25]{Interpolasi -- dan batasnya}
\begin{minipage}[t]{0.485\textwidth}
{\fontsize{9.4}{12.6}\selectfont\color{fg2}\textbf{Yang bisa dilakukan model}\par}\vskip 4pt
\begin{ibullets}
\item Memahami titik yang belum pernah dilihat dengan \textbf{mengaitkannya ke titik terdekat} di manifold.
\item Memahami keseluruhan ruang hanya dari sebagian sampelnya -- \textit{mengisi bagian yang kosong}.
\item Chollet menyebutnya \textbf{generalisasi setempat}.
\end{ibullets}

\end{minipage}%
\hfill
\begin{minipage}[t]{0.485\textwidth}
{\fontsize{9.4}{12.6}\selectfont\color{fg2}\textbf{Yang tidak bisa}\par}\vskip 4pt
\begin{ibullets}
\item \textbf{Generalisasi ekstrem} -- yang manusia lakukan setiap hari.
\item Anda bisa berpindah seminggu di NYC, seminggu di Shanghai, seminggu di Bangalore \hilite{tanpa ribuan kali seumur hidup berlatih untuk tiap kota}.
\item Itu ditopang abstraksi, model simbolik, penalaran, logika, akal sehat -- yang kita sebut \textbf{nalar}, bukan intuisi.
\end{ibullets}

\end{minipage}%
\begin{iband}[signal]
Peringatan Chollet: interpolasi hanya \hilite{puncak gunung es}. Menganggap interpolasi sama dengan seluruh generalisasi adalah kekeliruan. Bab 19 kembali ke sini.
\end{iband}
{\fontsize{9.4}{12.6}\selectfont\color{fg2}Perhatikan juga: interpolasi \textbf{pada manifold laten} berbeda dari interpolasi linear di ruang induk. Rerata piksel dua angka MNIST biasanya \hilite{bukan angka yang sah}.\par}\vskip 4pt
\end{frame}

\framekicker{Bagian 5.1.2}
\begin{frame}[shrink=25]{Mengapa deep learning bekerja -- dan syaratnya}
{\fontsize{9.4}{12.6}\selectfont\color{fg2}Model deep learning pada dasarnya adalah \textbf{kurva berdimensi sangat tinggi} yang mulus dan bersinambung (ia harus terdiferensialkan), dipasangkan ke titik data secara bertahap lewat penurunan gradien.\par}\vskip 4pt
\begin{isteps}
\item Kurva itu punya cukup parameter untuk memasangkan diri ke \textbf{apa saja} -- kalau dibiarkan cukup lama, ia murni menghafal.
\item Tetapi data yang dipasangkan \textbf{bukan titik-titik terpencar}; ia membentuk manifold berdimensi rendah yang sangat terstruktur.
\item Karena pemasangannya berlangsung \hilite{bertahap dan mulus}, ada titik di tengah pelatihan saat kurva model kira-kira menghampiri manifold alami data itu (gambar 5.10).
\end{isteps}
\begin{minipage}[t]{0.494\textwidth}
\begin{icardgood}
\cardh{Sifat 1 $\cdot$ pemetaan mulus dan bersinambung}
Wajib, karena harus terdiferensialkan. Kemulusan itu justru yang membantu menghampiri manifold laten.
\end{icardgood}
\end{minipage}%
\hfill
\begin{minipage}[t]{0.494\textwidth}
\begin{icardgood}
\cardh{Sifat 2 $\cdot$ prior arsitektur}
Strukturnya mencerminkan 'bentuk' informasi pada datanya -- terutama pada model citra (bab 8-12) dan model deret (bab 13).
\end{icardgood}
\end{minipage}%
\end{frame}

\framekicker{Bagian 5.1.3}
\begin{frame}[shrink=25]{Data latih adalah yang terpenting}
\begin{iquote}
\quotetext{Satu-satunya yang akan Anda temukan di dalam model deep learning adalah apa yang Anda masukkan ke dalamnya: prior yang tersandi di arsitekturnya, dan data yang dipakai melatihnya.}\quotecite{Chollet \& Watson, bab 5.1.3}
\end{iquote}
\begin{ibullets}
\item Kemampuan menggeneralisasi lebih merupakan akibat \textbf{struktur alami data} ketimbang sifat model.
\item Deep learning adalah pemasangan kurva, jadi ia butuh \textbf{pencuplikan yang padat} atas ruang masukannya -- terutama di dekat batas keputusan (gambar 5.11).
\item Cuplikan jarang $\rightarrow$ kurva yang dipelajari tidak cocok dengan ruang laten, dan \hilite{interpolasinya salah}.
\item Karena itu: \textbf{cara terbaik memperbaiki model adalah melatihnya dengan lebih banyak data, atau data yang lebih baik}.
\end{ibullets}
\begin{iband}[amber]
Kalau menambah data tidak mungkin, pilihan terbaik berikutnya adalah \hilite{membatasi banyaknya informasi yang boleh disimpan model}, atau menambah kekangan pada kemulusan kurvanya. Itulah \textbf{regularisasi}, dan itu isi bagian 5.4.4.
\end{iband}
\end{frame}

\coursesection{02}{Menilai model machine learning}{Anda hanya bisa mengendalikan yang bisa Anda amati.}

\framekicker{Bagian 5.2.1}
\begin{frame}[shrink=25]{Mengapa dua himpunan tidak cukup}
\begin{center}
\adjustbox{max width=\linewidth,max totalheight=0.52\textheight}{%

\begin{tikzpicture}[font=\sffamily\tiny]
  \node[draw=signal!60, fill=signal!16, rounded corners=3pt, minimum width=4.8cm,
        minimum height=0.5cm, text=ink] (tr) at (0,0) {latih --- model dipasangkan ke sini};
  \node[draw=amber!60, fill=amberbr!16, rounded corners=3pt, minimum width=2.3cm,
        minimum height=0.5cm, text=ink, anchor=west] (va) at (2.55,0) {validasi};
  \node[draw=lime!60, fill=limebr!16, rounded corners=3pt, minimum width=1.6cm,
        minimum height=0.5cm, text=ink, anchor=west] (te) at (5.0,0) {uji};
  \draw[amberbr, line width=0.9pt, dashed] (va.south) .. controls +(0,-0.5) and +(0,-0.5) .. (0.4,-0.25);
  \node[text=amber, anchor=west] at (-0.6,-0.95)
    {tiap kali hiperparameter disetel, sedikit informasi validasi BOCOR ke model};
  \node[draw=lime!45, fill=limebr!7, rounded corners=4pt, minimum width=9.0cm,
        minimum height=0.75cm, text=lime, align=left, anchor=north west] at (-2.4,-1.35)
    {~Data uji dipakai SEKALI, di akhir. Kalau ada satu saja setelan yang dipilih\\
     ~berdasarkan skor uji, ukuran generalisasi Anda sudah cacat.};
\end{tikzpicture}

}
\end{center}
\vskip 3pt{\color{ink3}\fontsize{7.4}{9.4}\selectfont Menyetel hiperparameter berdasarkan skor validasi adalah bentuk pembelajaran juga -- dan karenanya bisa overfit ke himpunan validasi.\par}
\begin{ibullets}
\item \textbf{Hiperparameter} (jumlah lapis, ukuran lapis) berbeda dari \textbf{parameter} (bobot). Yang pertama Anda setel; yang kedua dipelajari.
\item Tiap kali Anda menyetel satu hiperparameter berdasarkan skor validasi, \hilite{sedikit informasi tentang data validasi bocor ke model}.
\item Sekali dua kali tidak apa-apa. Diulang berkali-kali, himpunan validasi berhenti menjadi ukuran yang jujur.
\item Karena itu perlu himpunan ketiga: \textbf{data uji}, yang modelnya belum pernah menyentuhnya, bahkan tidak langsung.
\end{ibullets}
\end{frame}
\note{Untuk audit internal, ini bagian yang perlu dituliskan sebagai prosedur: siapa boleh melihat data uji, dan berapa kali.}

\framekicker{Bagian 5.2.1 $\cdot$ listing 5.5-5.6}
\begin{frame}[shrink=25]{Tiga protokol evaluasi}
\begin{itable}{>{\raggedright\arraybackslash}p{0.2068\textwidth}>{\raggedright\arraybackslash}p{0.3196\textwidth}>{\raggedright\arraybackslash}p{0.2444\textwidth}>{\raggedright\arraybackslash}p{0.1692\textwidth}}
\thead{Protokol} & \thead{Cara kerjanya} & \thead{Kapan dipakai} & \thead{Biayanya} \\ \midrule
\textbf{Hold-out sederhana} & Sisihkan sebagian sebagai uji; sisanya dibagi latih dan validasi. & Data banyak. & 1 model \\
\textbf{K-lipat} & Bagi jadi K bagian sama besar; tiap bagian sekali jadi validasi; skornya dirata-ratakan. & Skor sangat berayun bergantung belahan. & K model \\
\textbf{K-lipat berulang, diacak} & Jalankan K-lipat P kali, acak ulang tiap kali. & Data sedikit dan penilaian harus setepat mungkin. & P $\times$ K model \\
\end{itable}
\icode{python}{listing 5.6 --- inti K-lipat}{listings/ch05-04.py}
\begin{iband}[signal]
Perhatikan \inlinecode{model = get\_model()} di dalam gelung: tiap lipatan wajib memakai \hilite{instans yang benar-benar baru}. Memakai ulang model yang sudah terlatih membatalkan seluruh penilaian.
\end{iband}
\end{frame}

\framekicker{Bagian 5.2.2}
\begin{frame}[shrink=25]{Tolok banding akal sehat: altimeter roket yang tak terlihat}
\begin{iquote}
\quotetext{Melatih model deep learning itu seperti menekan tombol yang meluncurkan roket di dunia paralel. Anda tak bisa mendengarnya, tak bisa melihatnya. Satu-satunya umpan balik yang Anda punya adalah metrik validasi -- seperti altimeter pada roket yang tak terlihat itu.}\quotecite{Chollet \& Watson, bab 5.2.2}
\end{iquote}
\begin{itable}{>{\raggedright\arraybackslash}p{0.2444\textwidth}>{\raggedright\arraybackslash}p{0.2068\textwidth}>{\raggedright\arraybackslash}p{0.4888\textwidth}}
\thead{Persoalan} & \thead{Tolok banding akal sehat} & \thead{Alasannya} \\ \midrule
MNIST & \textbf{> 0,10} & Penebak acak atas 10 kelas seimbang. \\
IMDB & \textbf{> 0,50} & Dua kelas, seimbang 50:50. \\
Reuters & \textbf{$\approx$ 0,18-0,19} & 46 kelas, tetapi sebarannya tidak rata. \\
Biner 90:10 & \textbf{> 0,90} & Penebak yang selalu menjawab kelas A sudah dapat 0,90. Anda harus melampauinya. \\
\end{itable}
\begin{iband}[rose]
Kalau Anda \textbf{tidak bisa mengalahkan penyelesaian sepele}, model Anda tidak berharga. Mungkin modelnya salah, atau -- dan ini yang sering -- \hilite{persoalannya memang tidak bisa didekati dengan machine learning}. Kembali ke papan gambar.
\end{iband}
\end{frame}
\note{Slide ini yang paling sering menyelamatkan anggaran. Minta peserta menuliskan tolok banding untuk kasusnya sendiri SEBELUM melatih apa pun.}

\framekicker{Bagian 5.2.3}
\begin{frame}[shrink=25]{Tiga jebakan yang membatalkan penilaian}
\begin{minipage}[t]{0.3253\textwidth}
\begin{icardwarn}
\cardh{1 $\cdot$ Keterwakilan data}
Data terurut menurut kelas, lalu 80\% pertama diambil sebagai latih $\rightarrow$ latih hanya berisi kelas 0-7, uji hanya 8-9. \textit{Kelihatannya konyol, tetapi mengejutkan seringnya.} \textbf{Obatnya: acak sebelum membelah.}
\end{icardwarn}
\end{minipage}%
\hfill
\begin{minipage}[t]{0.3253\textwidth}
\begin{icardbad}
\cardh{2 $\cdot$ Anak panah waktu}
Meramal masa depan dari masa lalu? \textbf{Jangan diacak.} Mengacak menciptakan \hilite{kebocoran waktu}: model dilatih dengan data dari masa depan. Semua data uji harus lebih baru dari data latih.
\end{icardbad}
\end{minipage}%
\hfill
\begin{minipage}[t]{0.3253\textwidth}
\begin{icardbad}
\cardh{3 $\cdot$ Data kembar}
Titik data yang muncul dua kali -- lazim pada data dunia nyata. Setelah diacak, salinannya tersebar ke latih dan validasi: Anda menguji di atas data latih sendiri. \textbf{Yang terburuk dari semuanya.}
\end{icardbad}
\end{minipage}%
\begin{iband}[amber]
Pada data transaksional, jebakan \textbf{2 dan 3} hampir selalu muncul bersamaan: derenya bersifat temporal, dan subjek yang sama muncul berkali-kali. Membelah secara acak per baris \hilite{melanggar keduanya} dalam satu langkah.
\end{iband}
\end{frame}

\coursesection{03}{Memperbaiki fit}{Untuk mencapai fit yang pas, Anda harus overfit dulu.}

\framekicker{Bagian 5.3}
\begin{frame}[shrink=25]{Tiga kegagalan, tiga obat}
\begin{iquote}
\quotetext{Untuk mencapai fit yang sempurna, Anda harus overfit lebih dulu. Karena Anda tidak tahu di mana batasnya, Anda harus melewatinya untuk menemukannya.}\quotecite{Chollet \& Watson, bab 5.3}
\end{iquote}
\begin{itable}{>{\raggedright\arraybackslash}p{0.282\textwidth}>{\raggedright\arraybackslash}p{0.2632\textwidth}>{\raggedright\arraybackslash}p{0.3948\textwidth}}
\thead{Gejalanya} & \thead{Artinya} & \thead{Yang dicoba} \\ \midrule
\textbf{Pelatihan tidak mulai} -- rugi latih tidak turun & Setelan penurunan gradien salah. & Setel \textbf{learning rate} (turunkan atau naikkan) dan \textbf{ukuran batch}. Yang lain dibiarkan tetap. \\
\textbf{Mulai, tetapi tidak menggeneralisasi} -- tolok banding tak terkalahkan & Ada yang salah secara mendasar. & Mungkin datanya \textbf{tidak memuat informasi} untuk meramal targetnya; atau \textbf{prior arsitekturnya salah} untuk jenis data ini. \\
\textbf{Menggeneralisasi, tetapi tidak bisa overfit} & Kapasitas kurang. & Tambah lapis, perbesar lapis, atau pakai \textbf{jenis lapis yang lebih cocok}. \\
\end{itable}
\begin{iband}[rose]
Kegagalan kedua adalah \textbf{situasi terburuk} dalam machine learning: ia menandakan ada yang salah secara mendasar, dan \hilite{sering tidak mudah dikenali apanya}.
\end{iband}
\end{frame}

\framekicker{Bagian 5.3.1 $\cdot$ listing 5.7-5.8}
\begin{frame}[shrink=25]{Learning rate: satu angka yang menghentikan segalanya}
\begin{minipage}[t]{0.485\textwidth}
\icode{python}{listing 5.7 --- terlalu besar}{listings/ch05-05.py}
\ioutput{listings/ch05-06.txt}

\end{minipage}%
\hfill
\begin{minipage}[t]{0.485\textwidth}
\icode{python}{listing 5.8 --- wajar}{listings/ch05-07.py}
\ioutput{listings/ch05-08.txt}

\end{minipage}%
\begin{ibullets}
\item Learning rate \textbf{terlalu tinggi} $\rightarrow$ pembaruan melampaui jauh titik yang pas; hasilnya seperti di kiri.
\item \textbf{Terlalu rendah} $\rightarrow$ pelatihan begitu lambat sampai \hilite{tampak macet}, padahal sebetulnya jalan.
\item \textbf{Perbesar ukuran batch} $\rightarrow$ gradien lebih informatif dan kurang berisik (ragamnya lebih rendah).
\end{ibullets}
\begin{iband}[signal]
Semua parameter ini saling bergantung. Chollet menyarankan: cukup setel \hilite{learning rate dan ukuran batch saja}, sisanya biarkan tetap.
\end{iband}
\end{frame}

\framekicker{Bagian 5.3.3 $\cdot$ listing 5.9}
\begin{frame}[shrink=25]{Kapasitas: terlalu kecil, pas, terlalu besar}
\begin{itable}{>{\raggedright\arraybackslash}p{0.2068\textwidth}>{\raggedright\arraybackslash}p{0.282\textwidth}>{\raggedright\arraybackslash}p{0.3384\textwidth}>{\raggedright\arraybackslash}p{0.1128\textwidth}}
\thead{Model} & \thead{Arsitektur} & \thead{Yang terjadi} & \thead{Gambar} \\ \midrule
\textbf{Kapasitas kurang} & \inlinecode{Dense(10, softmax)} saja -- regresi logistik & Rugi validasi turun ke 0,26 lalu \textbf{diam di situ}. Bisa fit, tetapi \hilite{tidak bisa jelas-jelas overfit}. & 5.14 \\
\textbf{Kapasitas pas} & 2 $\times$ \inlinecode{Dense(128, relu)} & Fit cepat, mulai overfit \textbf{setelah 8 epoch}. Persis seperti seharusnya. & 5.15 \\
\textbf{Kapasitas berlebih} & 3 $\times$ \inlinecode{Dense(2048, relu)} & Overfit \textbf{langsung sejak awal}; rugi latih hampir nol dengan cepat, rugi validasi berisik. & 5.16 \\
\end{itable}
\icode{python}{listing 5.9 --- kapasitas kurang}{listings/ch05-09.py}
\begin{iband}[signal]
Alur kerjanya: \textbf{mulai dari sedikit lapis dan parameter, lalu besarkan sampai rugi validasi berhenti membaik}. Tidak ada rumus ajaib untuk menebak ukuran yang benar -- \hilite{harus dicoba, di himpunan validasi, bukan di himpunan uji}.
\end{iband}
\end{frame}

\coursesection{04}{Memperbaiki generalisasi}{Lima cara, berurut dari yang paling berdampak.}

\framekicker{Bagian 5.4.1}
\begin{frame}[shrink=25]{Kurasi data -- imbal hasil terbesar, dan sering dilewati}
\begin{iquote}
\quotetext{Deep learning adalah pemasangan kurva, bukan sihir.}\quotecite{Chollet \& Watson, bab 5.4.1}
\end{iquote}
\begin{iband}[signal]
Mengeluarkan lebih banyak tenaga dan uang untuk \textbf{pengumpulan data} hampir selalu memberi imbal hasil \hilite{jauh lebih besar} daripada jumlah yang sama untuk mengembangkan model yang lebih baik.
\end{iband}
\begin{isteps}
\item \textbf{Pastikan datanya cukup.} Anda butuh pencuplikan yang padat atas ruang masukan-silang-keluaran. Persoalan yang tampak mustahil kadang jadi bisa diselesaikan begitu datanya lebih banyak.
\item \textbf{Kecilkan galat pelabelan.} Lihat sendiri masukannya untuk mencari kejanggalan; periksa ulang labelnya.
\item \textbf{Bersihkan data dan tangani nilai yang hilang.} (Bab 6 membahasnya.)
\item \textbf{Lakukan seleksi fitur} kalau fiturnya banyak dan Anda belum yakin mana yang berguna.
\end{isteps}
\begin{iband}[amber]
Batasnya jelas: kalau persoalannya \textbf{terlalu berderau atau pada dasarnya diskret} -- misalnya mengurutkan senarai -- deep learning \hilite{tidak akan menolong}, seberapa pun datanya.
\end{iband}
\end{frame}

\framekicker{Bagian 5.4.2}
\begin{frame}[shrink=25]{Rekayasa fitur: contoh jam dinding}
\begin{itable}{>{\raggedright\arraybackslash}p{0.2444\textwidth}>{\raggedright\arraybackslash}p{0.3572\textwidth}>{\raggedright\arraybackslash}p{0.3384\textwidth}}
\thead{Tingkat} & \thead{Masukannya} & \thead{Yang dibutuhkan} \\ \midrule
\textbf{Data mentah} & Kisi piksel citra jam & ConvNet, dan sumber daya komputasi yang tidak sedikit. \\
\textbf{Fitur lebih baik} & Koordinat (x, y) ujung tiap jarum & Algoritma machine learning sederhana sudah cukup. \\
\textbf{Lebih baik lagi} & Sudut $\theta$ tiap jarum (koordinat polar) & \hilite{Tidak perlu machine learning sama sekali} -- cukup pembulatan dan pencarian di kamus. \\
\end{itable}
{\fontsize{9.4}{12.6}\selectfont\color{fg2}Itulah inti rekayasa fitur: \textbf{membuat persoalan lebih mudah dengan menyatakannya secara lebih sederhana}. Membuat manifold latennya lebih mulus, lebih sederhana, lebih rapi. Ia menuntut pemahaman mendalam atas persoalannya.\par}\vskip 4pt
\begin{minipage}[t]{0.494\textwidth}
\begin{icardgood}
\cardh{Tetap berguna --- alasan 1}
Fitur yang baik menyelesaikan persoalan \textbf{dengan sumber daya lebih sedikit}. Memakai ConvNet untuk membaca jam dinding itu konyol.
\end{icardgood}
\end{minipage}%
\hfill
\begin{minipage}[t]{0.494\textwidth}
\begin{icardgood}
\cardh{Tetap berguna --- alasan 2}
Fitur yang baik menyelesaikan persoalan dengan \textbf{data jauh lebih sedikit}. Kemampuan model menemukan fiturnya sendiri bergantung pada tersedianya banyak data latih.
\end{icardgood}
\end{minipage}%
\begin{iband}[signal]
Sebelum deep learning, rekayasa fitur adalah \textbf{bagian terpenting} alur kerja machine learning. Solusi MNIST dulu dibangun dari fitur yang ditulis tangan: jumlah gelung pada citra angka, tinggi tiap angka, histogram nilai piksel.
\end{iband}
\end{frame}

\framekicker{Bagian 5.4.3}
\begin{frame}[shrink=25]{Early stopping}
{\fontsize{9.4}{12.6}\selectfont\color{fg2}Model deep learning \textbf{selalu berparameter berlebihan}: derajat kebebasannya jauh melebihi yang minimal diperlukan. Itu bukan masalah, sebab Anda \hilite{tidak pernah memasangkannya sampai penuh}. Pelatihan selalu dihentikan jauh sebelum rugi latih minimum tercapai.\par}\vskip 4pt
\begin{minipage}[t]{0.485\textwidth}
{\fontsize{9.4}{12.6}\selectfont\color{fg2}\textbf{Cara di bab 4} -- latih lebih lama dari perlu untuk menemukan epoch terbaik, lalu latih model baru tepat sebanyak epoch itu.\par}\vskip 4pt
\begin{iband}[amber]
Lazim, tetapi menuntut \hilite{pekerjaan berulang} yang kadang mahal.
\end{iband}

\end{minipage}%
\hfill
\begin{minipage}[t]{0.485\textwidth}
{\fontsize{9.4}{12.6}\selectfont\color{fg2}\textbf{Cara yang lebih baik} -- simpan model tiap akhir epoch, lalu pakai yang terbaik. Di Keras ini dikerjakan callback \inlinecode{EarlyStopping}, yang menghentikan pelatihan begitu metrik validasi berhenti membaik \textbf{sambil mengingat keadaan model terbaik}.\par}\vskip 4pt
\begin{iband}[signal]
Callback dibahas penuh di \textbf{bab 7}.
\end{iband}

\end{minipage}%
{\fontsize{9.4}{12.6}\selectfont\color{fg2}Menemukan titik persis saat fit paling bisa digeneralisasi -- batas antara kurva underfit dan overfit -- adalah \textbf{salah satu hal paling ampuh} yang bisa Anda lakukan untuk memperbaiki generalisasi.\par}\vskip 4pt
\end{frame}

\framekicker{Bagian 5.4.4 $\cdot$ listing 5.10-5.12}
\begin{frame}[shrink=25]{Regularisasi 1: kecilkan modelnya}
{\fontsize{9.4}{12.6}\selectfont\color{fg2}Kalau model punya sumber daya penghafalan yang terbatas, ia \hilite{terpaksa mempelajari representasi terpadat yang punya daya ramal} -- persis jenis representasi yang kita inginkan.\par}\vskip 4pt
\begin{itable}{>{\raggedright\arraybackslash}p{0.1692\textwidth}>{\raggedright\arraybackslash}p{0.188\textwidth}>{\raggedright\arraybackslash}p{0.1504\textwidth}>{\raggedright\arraybackslash}p{0.3196\textwidth}>{\raggedright\arraybackslash}p{0.1128\textwidth}}
\thead{Versi} & \thead{Lapis antara} & \thead{Mulai overfit} & \thead{Perilakunya} & \thead{Gambar} \\ \midrule
\textbf{Acuan} & 2 $\times$ \inlinecode{Dense(16)} & epoch 4 & Titik banding. & --- \\
\textbf{Lebih kecil} & 2 $\times$ \inlinecode{Dense(4)} & epoch 6 & Mulai overfit lebih lambat, dan \textbf{memburuknya lebih pelan}. & 5.18 \\
\textbf{Jauh lebih besar} & 2 $\times$ \inlinecode{Dense(512)} & epoch 1 & Overfit hampir seketika dan jauh lebih parah; rugi validasinya juga lebih berisik. & 5.19 \\
\end{itable}
\begin{iband}[amber]
Tidak ada rumus ajaib untuk ukuran yang benar. Alur kerjanya: \textbf{mulai dari sedikit, besarkan sampai imbal hasil rugi validasi mengecil} -- dan lakukan itu di \hilite{himpunan validasi, tentu saja, bukan himpunan uji}.
\end{iband}
\end{frame}

\framekicker{Bagian 5.4.4 $\cdot$ listing 5.13-5.14}
\begin{frame}[shrink=25]{Regularisasi 2: kekang bobotnya (L1 / L2)}
\begin{iquote}
\quotetext{Pisau cukur Occam: diberi dua penjelasan atas suatu hal, yang paling mungkin benar adalah yang paling sederhana -- yang paling sedikit andaiannya.}\quotecite{Prinsip yang dipakai Chollet untuk membenarkan regularisasi bobot}
\end{iquote}
\icode{python}{listing 5.13-5.14}{listings/ch05-10.py}
\begin{itable}{>{\raggedright\arraybackslash}p{0.1504\textwidth}>{\raggedright\arraybackslash}p{0.4324\textwidth}>{\raggedright\arraybackslash}p{0.3572\textwidth}}
\thead{Jenis} & \thead{Biaya yang ditambahkan ke rugi} & \thead{Efeknya} \\ \midrule
\textbf{L1} & Sebanding dengan \textbf{nilai mutlak} koefisien bobot (norma L1). & Mendorong bobot menjadi \textbf{jarang} (banyak yang nol). \\
\textbf{L2} & Sebanding dengan \textbf{kuadrat} koefisien bobot (norma L2). & Mendorong semua bobot menjadi kecil. Disebut juga \textbf{weight decay} -- namanya beda, matematikanya sama. \\
\end{itable}
\begin{iband}[signal]
\inlinecode{l2(0.002)} berarti tiap koefisien menambahkan \inlinecode{0.002 * nilai\_bobot ** 2} ke rugi total. Denda ini hanya ditambahkan \textbf{saat latih}, jadi \hilite{rugi saat latih akan tampak jauh lebih tinggi daripada saat uji} -- dan itu normal, bukan bug.
\end{iband}
\end{frame}

\framekicker{Bagian 5.4.4 $\cdot$ listing 5.15}
\begin{frame}[shrink=25]{Regularisasi 3: dropout}
\begin{center}
\adjustbox{max width=\linewidth,max totalheight=0.52\textheight}{%

\begin{tikzpicture}[font=\sffamily\tiny]
  \node[draw=rule, fill=papertint, rounded corners=3pt, minimum width=2.6cm,
        minimum height=1.5cm, text=ink, font=\ttfamily\tiny, align=left] (a) at (0,0)
    {0.3~~0.2~~1.5~~0.0\\0.6~~0.1~~0.0~~0.3\\0.2~~1.9~~0.3~~1.2\\0.7~~0.5~~1.0~~0.0};
  \node[text=ink3, anchor=south] at (a.north) {keluaran lapis};
  \node[text=amber] at (1.85,0.2) {50\% dropout};
  \draw[amberbr, line width=0.9pt, -{Stealth[length=4pt]}] (1.4,-0.15) -- (2.35,-0.15);
  \node[draw=signal!60, fill=signal!9, rounded corners=3pt, minimum width=2.6cm,
        minimum height=1.5cm, text=ink, font=\ttfamily\tiny, align=left] (b) at (3.9,0)
    {\textcolor{rose}{0.0}~~0.2~~1.5~~\textcolor{rose}{0.0}\\
     0.6~~0.1~~\textcolor{rose}{0.0}~~0.3\\
     \textcolor{rose}{0.0}~~1.9~~0.3~~\textcolor{rose}{0.0}\\
     0.7~~\textcolor{rose}{0.0}~~\textcolor{rose}{0.0}~~\textcolor{rose}{0.0}};
  \node[text=lime, font=\small] at (5.5,0) {$\times 2$};
  \node[draw=lime!50, fill=limebr!10, rounded corners=3pt, minimum width=3.3cm,
        minimum height=1.5cm, text=ink2, align=left, anchor=west] at (6.1,0)
    {diskalakan naik saat LATIH,\\supaya saat UJI matriksnya\\tidak perlu diubah sama sekali\\[2pt]
     \textcolor{lime}{\ttfamily layer\_output /= 0.5}};
  \node[text=amber, anchor=west] at (-1.5,-1.15)
    {Gambar 5.21 --- laju dropout lazimnya 0,2 sampai 0,5. Saat uji tidak ada unit yang dijatuhkan.};
\end{tikzpicture}

}
\end{center}
\vskip 3pt{\color{ink3}\fontsize{7.4}{9.4}\selectfont Gambar 5.21 -- dropout pada matriks aktivasi saat latih, dengan penskalaan dikerjakan saat latih juga.\par}
\icode{python}{listing 5.15 --- dropout pada model IMDB}{listings/ch05-11.py}
\end{frame}

\framekicker{Bagian 5.4.4}
\begin{frame}[shrink=25]{Dari mana ide dropout datang -- kisah Hinton dan bank}
\begin{iquote}
\quotetext{Saya pergi ke bank saya. Tellernya berganti terus dan saya tanya salah satunya kenapa. Dia bilang tidak tahu, tetapi mereka memang sering dipindah-pindah. Saya kira itu pasti karena menipu bank memerlukan kerja sama antar-pegawai. Dari situ saya sadar bahwa membuang subhimpunan neuron yang berbeda secara acak pada tiap contoh akan mencegah persekongkolan, dan dengan begitu mengurangi overfitting.}\quotecite{Geoff Hinton, dikutip Chollet \& Watson, bab 5.4.4}
\end{iquote}
\begin{iband}[signal]
Gagasan intinya: menyuntikkan derau ke nilai keluaran sebuah lapis \hilite{memutus pola kebetulan yang tidak bermakna} -- yang Hinton sebut \textit{persekongkolan} -- yang akan mulai dihafal model bila tidak ada derau sama sekali.
\end{iband}
\begin{iband}[amber]
Asal-usulnya kebetulan menolong sebagai analogi: \textbf{rotasi pegawai sebagai kendali kecurangan} adalah kendali internal yang sudah dikenal luas. Dropout adalah kendali yang sama, diterapkan pada neuron.
\end{iband}
{\fontsize{9.4}{12.6}\selectfont\color{fg2}Hasilnya pada model IMDB (gambar 5.22): peningkatan yang jelas atas model acuan, dan \textbf{tampak bekerja lebih baik daripada regularisasi L2}, sebab rugi validasi terendah yang dicapainya lebih rendah.\par}\vskip 4pt
\end{frame}

\framekicker{Bagian 5.4.4}
\begin{frame}[shrink=25]{Mana yang dipakai kapan}
\begin{itable}{>{\raggedright\arraybackslash}p{0.2256\textwidth}>{\raggedright\arraybackslash}p{0.3572\textwidth}>{\raggedright\arraybackslash}p{0.3572\textwidth}}
\thead{Cara} & \thead{Kapan paling cocok} & \thead{Catatan} \\ \midrule
\textbf{Data lebih banyak / lebih baik} & Selalu, kalau mungkin. & Imbal hasil terbesar. Menambah data yang terlalu berderau justru merugikan. \\
\textbf{Fitur yang lebih baik} & Data sedikit; persoalan yang dipahami mendalam. & Bisa menghapus kebutuhan akan model besar sama sekali. \\
\textbf{Kurangi kapasitas} & Model kecil; data sedikit. & Cari kompromi -- jangan sampai malah underfit. \\
\textbf{Regularisasi bobot (L1/L2)} & \textbf{Model deep learning yang lebih kecil}. & Pada model besar yang sangat berparameter lebih, mengekang nilai bobot \hilite{tidak banyak berpengaruh}. \\
\textbf{Dropout} & \textbf{Model besar} -- di situ regularisasi bobot kurang mempan. & Laju lazim 0,2-0,5. Pada IMDB hasilnya mengalahkan L2. \\
\end{itable}
\begin{iband}[rose]
Satu syarat menaungi semuanya: \textbf{regularisasi harus selalu dipandu prosedur evaluasi yang tepat}. Anda hanya akan mencapai generalisasi \hilite{kalau Anda bisa mengukurnya}.
\end{iband}
\end{frame}

\framekicker{Ringkasan}
\begin{frame}[shrink=25]{Yang wajib terbawa dari bab 5}
\begin{isteps}
\item Tujuan model adalah \textbf{menggeneralisasi}: tepat pada masukan yang belum pernah dilihat. Lebih sulit daripada kelihatannya.
\item Jaringan dalam menggeneralisasi dengan \textbf{menginterpolasi} pada manifold laten data. Karena itu ia hanya memahami masukan yang \hilite{dekat dengan yang pernah dilihatnya}.
\item Persoalan pokoknya adalah \textbf{ketegangan optimalisasi vs generalisasi}. Setiap praktik baik di buku ini menangani ketegangan itu.
\item \textbf{Ukur dulu, baru perbaiki.} Hold-out, K-lipat, K-lipat berulang -- dan selalu sisakan himpunan uji yang benar-benar tak tersentuh.
\item \textbf{Tetapkan tolok banding akal sehat} sebelum melatih. Kalau tak terkalahkan, mungkin persoalannya bukan persoalan machine learning.
\item Untuk fit: setel \textbf{learning rate dan ukuran batch}, perbaiki \textbf{prior arsitektur}, tambah \textbf{kapasitas} sampai bisa overfit.
\item Untuk generalisasi: \textbf{data lebih baik $\rightarrow$ fitur lebih baik $\rightarrow$ early stopping $\rightarrow$ kurangi kapasitas $\rightarrow$ L1/L2 $\rightarrow$ dropout}, dalam urutan itu.
\end{isteps}
\vskip 4pt
\reslink{NOTEBOOK}{02\_label\_diacak\_manifold.ipynb}{../../course-slides/notebooks/ch05/02\_label\_diacak\_manifold.ipynb}
\reslink{NOTEBOOK}{05\_regularisasi\_l2\_dropout.ipynb}{../../course-slides/notebooks/ch05/05\_regularisasi\_l2\_dropout.ipynb}
\reslink{BAB BERIKUT}{Bab 6 --- Alur kerja universal}{../ch06/index.html}
\vskip 2pt
\end{frame}

\end{document}
